\chapter*{Introduction}
\addcontentsline{toc}{chapter}{Introduction}
\markboth{Introduction}{}
% BUDGET : 3 pages.

% 1. Le territoire et sa dependance alimentaire : ce que TI-FIG cherche a deplacer.
\arediger{la question alimentaire guadeloupéenne, et le projet TI-FIG}

% 2. La commande : concevoir des scenarios agro-economiques de developpement des
%    aliments traditionnels, avec MOSAICA comme outil ATTENDU.
\arediger{la commande telle qu'elle a été formulée}

% 3. Le constat qui a reoriente le stage -- le coeur du memoire. Le modele n'etait
%    pas exploitable en l'etat pour produire les scenarios demandes.
\arediger{pourquoi il a fallu reconstruire avant de pouvoir simuler}

% 4. Positionnement : d'ou je parle. Ingenieur en mecanique des fluides et
%    energetique, six mois, pas d'enquete de terrain, des donnees produites par
%    d'autres. Ce que cette position permet de voir et ce qu'elle interdit.
\arediger{section de positionnement}

% 5. Problematique et annonce du plan.
\arediger{problématique et annonce du plan}
